\chapter{Spider silk, from first principles}
\label{ch:biology}

\section{What a protein is}

A protein is a chain. Each link in the chain is an \term{amino acid}, and there
are twenty different kinds that living things routinely use. Chemistry aside,
the only thing that matters for these notes is this: a protein can be written
down as a string of letters, one letter per link, drawn from a
twenty-letter alphabet.

\begin{center}
\aaseq{M K H T I G I L G G M G P A A T A D M L E K \dots}
\end{center}

The twenty letters, with the names worth knowing for silk:

\begin{center}
\small
\begin{tabular}{llllll}
\toprule
A alanine & C cysteine & D aspartate & E glutamate & F phenylalanine & G glycine \\
H histidine & I isoleucine & K lysine & L leucine & M methionine & N asparagine \\
P proline & Q glutamine & R arginine & S serine & T threonine & V valine \\
W tryptophan & Y tyrosine & & & & \\
\bottomrule
\end{tabular}
\end{center}

A chain of a few hundred of these folds up into a three-dimensional shape, and
the shape plus the chemistry determines what the protein does. For silk, the
relevant consequence is mechanical: how the chains pack together determines how
strong and how stretchy the resulting fibre is.

\begin{keybox}{The central premise of the whole field}
The sequence determines the structure, and the structure determines the
properties. So in principle, if you could learn the map from
\emph{letters} to \emph{mechanical properties}, you could design a fibre by
choosing letters. That map is what the paper tries to learn. Whether it
succeeded, and in what sense, is what this project is about.
\end{keybox}

\section{Why spider silk specifically}

Spider dragline silk --- the stuff a spider hangs from, and the radial spokes
of a web --- has an unusual combination of properties. It is roughly as strong
as steel by weight, far stretchier, and therefore absorbs much more energy
before breaking. That last quantity, energy absorbed per unit volume, is called
\term{toughness}, and it is the property spider silk is famous for.

Four mechanical properties recur throughout this project. They are all read off
a single experiment: take a fibre, pull it until it snaps, and record force
against extension.

\begin{figure}[htbp]
\centering
\begin{tikzpicture}[font=\small]
\begin{axis}[
  width=11cm, height=6.4cm,
  xlabel={strain (extension, \%)},
  ylabel={stress (force per area)},
  xmin=0,xmax=32,ymin=0,ymax=1.25,
  axis lines=left, xtick=\empty, ytick=\empty,
  every axis plot/.append style={thick},
  clip=false,
]
\addplot[papercol,line width=1.2pt,domain=0:28,samples=120]
  {x<6 ? 0.115*x : 0.69 + 0.30*(1-exp(-(x-6)/7)) + 0.010*(x-6)};
\addplot[ourscol,dashed,line width=0.8pt] coordinates {(0,0) (6,0.69)};

% breaking point
\node[circle,fill=ourscol,inner sep=1.6pt] at (axis cs:28,1.155) {};
\node[anchor=west,font=\scriptsize,text=ourscol] at (axis cs:28.4,1.155)
  {breaks here};

% annotations
\node[anchor=west,font=\scriptsize] at (axis cs:1.2,0.86)
  {\textbf{slope} $=$ elastic modulus $E$};
\draw[->,greyc,thin] (axis cs:4.6,0.83) -- (axis cs:3.3,0.42);

\node[anchor=east,font=\scriptsize] at (axis cs:27,0.62)
  {\textbf{height at break} $=$ tensile strength};
\draw[->,greyc,thin] (axis cs:27,0.68) -- (axis cs:27.8,1.10);

\node[anchor=north,font=\scriptsize] at (axis cs:16,0.30)
  {\textbf{area under curve} $=$ toughness};

\node[anchor=north,font=\scriptsize] at (axis cs:28,0.06)
  {\textbf{strain at break}};

\addplot[fill=papercol,opacity=0.10,draw=none,domain=0:28,samples=120]
  {x<6 ? 0.115*x : 0.69 + 0.30*(1-exp(-(x-6)/7)) + 0.010*(x-6)} \closedcycle;
\end{axis}
\end{tikzpicture}
\caption[Stress--strain curve]{A stress--strain curve, and where the four
mechanical properties come from. Pull a fibre until it snaps. The initial
\emph{slope} is stiffness (elastic modulus $E$); the \emph{height} at failure
is tensile strength; the \emph{horizontal position} of failure is strain at
break; the \emph{shaded area} is toughness, the energy the fibre absorbed. All
four come from one test, which is why they appear together in the dataset.}
\label{fig:stressstrain}
\end{figure}

\begin{keybox}{The four properties, and why there are eight numbers}
\begin{description}[leftmargin=2.6cm,style=nextline,itemsep=2pt]
  \item[Toughness] Energy absorbed per unit volume before breaking. The area
        under the curve. Units GJ/m$^3$.
  \item[Elastic modulus $E$] Stiffness: how much force to stretch it a little.
        The initial slope. Units GPa.
  \item[Tensile strength] The maximum stress it withstands. Units GPa.
  \item[Strain at break] How far it stretches before failing, as a percentage.
\end{description}
Each is measured several times on fibres from the same spider, so each also has
a \term{standard deviation} --- a measure of how much the repeat measurements
disagreed. Four properties plus four standard deviations makes the
\term{eight-dimensional property vector} that runs through this entire project.
\end{keybox}

The standard deviations are not decoration. A spider's dragline is not made of
one protein but of several, in proportions nobody fully knows, so the same
sequence can be associated with fibres that test differently. Including the
spread is the authors' way of carrying that uncertainty into the model.

\section{Spidroins, and the hierarchy that makes silk work}

The proteins in spider silk are called \term{spidroins}. A spider makes several
families of them from different glands, for different jobs:

\begin{center}
\small
\begin{tabular}{lll}
\toprule
\textbf{Family} & \textbf{Full name} & \textbf{Job} \\
\midrule
\textbf{MaSp} & major ampullate spidroin & dragline; web spokes. \emph{The subject of this paper.} \\
MiSp & minor ampullate spidroin & auxiliary spiral \\
Flag & flagelliform spidroin & the stretchy capture spiral \\
AcSp & aciniform spidroin & prey wrapping, egg-sac lining \\
AgSp & aggregate spidroin & the sticky glue coating \\
PySp & pyriform spidroin & attachment cement \\
CySp & cylindriform (tubuliform) & egg-sac outer wall \\
\bottomrule
\end{tabular}
\end{center}

MaSp is the one that matters here, because dragline silk is what gets
mechanically tested. The paper trains on MaSp only, and one of my extensions
(Section~\ref{sec:nonmasp}) asks what the model does when handed the others.

A spidroin has a characteristic layout, and the layout explains the mechanics.

\begin{figure}[htbp]
\centering
\begin{tikzpicture}[font=\small,
  dom/.style={draw=greyc!70,rounded corners=2pt,minimum height=0.85cm},
  rep/.style={dom,fill=papercol!18,draw=papercol!60},
  nt/.style={dom,fill=ourscol!18,draw=ourscol!60},
]
% the chain
\node[nt,minimum width=1.9cm] (n) at (0,0) {N-terminal};
\node[rep,minimum width=1.35cm,right=1pt of n] (r1) {poly-A};
\node[rep,minimum width=1.55cm,right=1pt of r1] (r2) {GGX};
\node[rep,minimum width=1.35cm,right=1pt of r2] (r3) {poly-A};
\node[rep,minimum width=1.55cm,right=1pt of r3] (r4) {GPGXX};
\node[rep,minimum width=1.35cm,right=1pt of r4] (r5) {poly-A};
\node[nt,minimum width=1.9cm,right=1pt of r5] (c) {C-terminal};

\draw[decorate,decoration={brace,amplitude=5pt,mirror},greyc]
  ($(r1.south west)+(0,-3pt)$) -- ($(r5.south east)+(0,-3pt)$)
  node[midway,below=7pt,font=\scriptsize,text width=6cm,align=center]
  {\textbf{repetitive core}: the same short motifs, over and over, sometimes
   for thousands of residues};

\node[above=3pt of n,font=\scriptsize,text width=2.2cm,align=center,text=ourscol]
  {solubility,\\ assembly};
\node[above=3pt of c,font=\scriptsize,text width=2.2cm,align=center,text=ourscol]
  {solubility,\\ assembly};

% zoom: what polyA and GGX do
\node[draw=greyc!50,fill=boxbg,rounded corners=2pt,inner sep=6pt,
      text width=5.2cm,align=left,font=\scriptsize,below=1.9cm of r2.south]
  (z1) {\textbf{poly-A} $=$ \aaseq{AAAAAAA}\\
        Chains line up side by side into rigid, crystalline
        \emph{$\beta$-sheet} stacks. These are the hard reinforcing
        blocks. They supply \emph{strength} and \emph{stiffness}.};
\node[draw=greyc!50,fill=boxbg,rounded corners=2pt,inner sep=6pt,
      text width=5.2cm,align=left,font=\scriptsize,below=1.9cm of r4.south]
  (z2) {\textbf{GGX} and \textbf{GPGXX}\\
        Floppy, disordered, spring-like regions. They uncoil under load.
        These supply \emph{extensibility}, and therefore most of the
        \emph{toughness}.};
\draw[greyc,thin,dashed] (r1.south) -- (z1.north);
\draw[greyc,thin,dashed] (r4.south) -- (z2.north);
\end{tikzpicture}
\caption[Spidroin layout]{The layout of a spidroin. Two non-repetitive terminal
domains bracket a long repetitive core. The core alternates hard crystalline
blocks with soft extensible ones. ``X'' stands for a position that varies.}
\label{fig:spidroin}
\end{figure}

\begin{keybox}{The analogy that makes silk make sense}
Think of a rope made of rubber bands with steel beads threaded on at intervals.
The steel beads (poly-alanine $\beta$-sheets) stop it from pulling apart --- they
give strength. The rubber (GGX, GPGXX) lets it stretch a long way before
anything breaks --- it gives extensibility. Toughness, the energy absorbed, is
strength \emph{times} extensibility, roughly, which is why you need both, and
why silk beats materials that have only one.
\end{keybox}

This is why the ablation experiment in Section~\ref{sec:ablations} deletes
poly-A tracts specifically. If the model has learned anything mechanically
meaningful, removing the hard blocks should change its predictions about
strength and stiffness. Whether it does is one of the project's main findings.

\section{The Silkome dataset}
\label{sec:silkome}

Everything downstream depends on one dataset, so it is worth being precise
about what is in it.

\begin{keybox}{Silkome}
Arakawa \emph{et al.}, ``1000 spider silkomes: linking sequences to silk
physical properties'', \emph{Science Advances} \textbf{8}, eabo6043 (2022).

Two files matter:
\begin{itemize}[leftmargin=*,itemsep=2pt]
  \item \code{spider-silkome-database.v1.prot.fasta} --- 11{,}155 spidroin
        sequence records across many species.
  \item \code{mechanical\_properties.csv} --- 446 rows, one per individual
        spider whose dragline was pulled apart in a testing machine.
\end{itemize}
\end{keybox}

Two features of this dataset caused me real work later.

\paragraph{The sequences are terminal-domain records, not whole proteins.}
Each FASTA record is labelled \code{NTD} or \code{CTD}, meaning it is anchored
on the N- or C-terminal domain and runs some distance into the repetitive core.
A full spidroin can be several thousand residues; these records are typically
a few hundred. So when the paper says it trains on ``sequences'', it means
these fragments.

\paragraph{Properties belong to spiders, sequences belong to genes.}
A single spider yields many sequence records --- MaSp1, MaSp2, MiSp, Flag and so
on, each NTD and CTD --- but only \emph{one} set of mechanical measurements,
taken on its dragline. So the join between the two files is one-to-many, and
the key is the individual spider. Getting this wrong is what made my first
attempt at rebuilding the dataset return zero rows
(Section~\ref{sec:dataset}).

\begin{warnbox}{A consequence worth carrying forward}
Because the mechanical test is done on \emph{dragline}, and dragline is made of
\emph{MaSp}, a Flag sequence in this dataset is paired with the mechanical
properties of the same spider's dragline --- not of its flagelliform silk. The
paper is explicit about assuming ``MaSp governs the mechanics'', and restricting
to MaSp follows from that. It also means the out-of-distribution experiment in
Section~\ref{sec:nonmasp} has to be phrased carefully: no analysis of this
dataset could tell you whether a model predicts flagelliform properties,
because the dataset does not contain them.
\end{warnbox}
